\documentclass[10pt,a4paper]{amsart}
\usepackage[margin=19mm]{geometry}
\usepackage{amsmath,amssymb,mathtools}
\usepackage[scr=rsfs,cal=euler]{mathalfa}
\usepackage{xfrac}
\usepackage[unicode,hidelinks,bookmarks=false]{hyperref}
\newcommand{\ie}{i.e.\ }
\newcommand{\cf}{cf.\ }
\newcommand{\resp}{respectively\ }
\newcommand{\Subsection}{\subsection}
\providecommand{\nobreakdash}{\nobreak\hbox{-}}
\setlength{\emergencystretch}{2em}
\multlinegap0pt
\usepackage{amssymb}
\usepackage[shortlabels]{enumitem}
\newtheorem{lemma}{Lemma}
\newtheorem{proposition}{Proposition}
\newcommand{\on}{\operatorname}
\newcommand{\ord}{\on{ord}}
\title{The unramified computation of a Shimura integral for $\mathrm{SL}(2)\times \mathrm{GL}(2)$}
\author{Pan Yan}
\date{Source: arXiv:2210.13653v1 (CC BY 4.0)}
\begin{document}
\maketitle

\noindent\emph{Source and licence.} The unramified computation of a Shimura integral for $\mathrm{SL}(2)\times \mathrm{GL}(2)$. Version arXiv:2210.13653v1 (CC BY 4.0), distributed by arXiv under the Creative Commons Attribution 4.0 International licence, http://creativecommons.org/licenses/by/4.0/, as recorded in the arXiv licence metadata for this exact version and shown on the abstract page https://arxiv.org/abs/2210.13653v1. Full licence notice: \texttt{licenses/arxiv-2210.13653-CC-BY-4.0.txt}.\par\smallskip

\noindent\textit{Editorial scope.} Counted unit: the article's proposition eq-unramified-J2 (source0.tex lines 444--451). The article's complete original proof of it is delivered with it (source0.tex lines 454--523, one \texttt{proof} environment of 70 lines). No mathematics is written, repaired or re-derived: the statement and the proof are the article's own lines, in the article's own notation, and no retained line is substituted or re-flowed.  Retained uncounted as the source-local context the proof needs: lines 326--344; lines 525--531.  Nothing was omitted from any retained span, so the package carries no omission record.  The retained lines cite 2 works, whose entries are transcribed verbatim from the article's own bibliography source in the e-print and delivered with short editorial labels recorded in the item's reference mapping.  No retained line contains a figure environment, an image-inclusion command or drawing code, so no image is delivered and the package declares no asset dependency.  Editorial formatting only, in addition: an AMS \texttt{amsart} preamble that restates the article's own declarations and macros used by the retained lines, namely the environments \texttt{lemma}, \texttt{proposition} on separate counters, together with the standard packages those lines need.  The retained environments print in this document as Lemma 1, Proposition 1 in order of appearance, whereas the article prints the same items with the numbering of its own full document.  The complete native source of the article as distributed in the arXiv e-print for this exact version is bundled unchanged as \texttt{sources/132547/source0.tex}.
\begin{lemma}\cite[Theorem 5.4]{CasselmanShalika1980II}
Let $a\in F^\times$. Then 
\begin{equation*}
W_\pi^0  \left( \begin{smallmatrix} a & \\ & a^{-1}\end{smallmatrix}\right) =  
 \begin{cases}
|a| \cdot \frac{\chi(\varpi)^{\ord(a)+1}-\chi(\varpi)^{-\ord(a)}}{\chi(\varpi)-1}   &\text{ if } \ord(a)\ge 0 \\
0 & \text{ if }\ord(a)< 0
\end{cases}
\end{equation*}
and 
\begin{equation*}
W_\tau^0\left( \begin{smallmatrix} a & \\&1\end{smallmatrix}\right)=
 \begin{cases}
 |a|^{\frac{1}{2}} \cdot   \frac{\chi_1(\varpi)^{\ord(a)+1}-\chi_2(\varpi)^{\ord(a)+1} }{\chi_1(\varpi)-\chi_2(\varpi)} &\text{ if }\ord(a)\ge 0\\
 0 &\text{ if }\ord(a)<0
 \end{cases}.
 \end{equation*}  
 \label{lemma-CasselmanShalika}
\end{lemma}

\begin{lemma}
For any $m\ge 2$, we have
$$
  \int\limits_{\mathcal{O}_F^\times} \gamma_\psi( \varpi^{m} u^{-1})^{-1}  \psi( \varpi^{-m} u)du=0.
$$
\label{lemma-local-integral-gamma-psi-vanishing}
\end{lemma}

\begin{proposition}
For $a\in F^\times \cap \mathcal{O}_F$, we have 
\begin{equation}
J_2(a)=   |a|^{\frac{1}{2}} \cdot \frac{\chi_1(\varpi)^{\ord(a)}-\chi_2(\varpi)^{\ord(a)}}{\chi_1(\varpi)-\chi_2(\varpi)}\cdot \chi_1(\varpi)\chi_2(\varpi) \cdot q^{-(s+1/2)}.
\label{eq-unramified-J2}
\end{equation}
\label{prop-unramified-J2}
\end{proposition}

\begin{proof}
To deal with the integral when $z\in F\setminus \mathcal{O}_F$, we consider the following matrix identity
\begin{equation*}
\begin{split}
\gamma  u(0,0,z) \gamma^{-1}   &=   \left( \begin{smallmatrix} 1 &&&\\&1&&\\&-z&1&\\ &&&1 \end{smallmatrix}\right) =     \left(\begin{smallmatrix} 1 &&&\\&z^{-1}&-1&\\&&z&\\ &&&1\end{smallmatrix}\right)    \left(\begin{smallmatrix} 1 &&&\\&&1&\\&-1&&\\ &&&1\end{smallmatrix}\right)   \left(\begin{smallmatrix} 1 &&&\\&1&-z^{-1}&\\&&1&\\ &&&1\end{smallmatrix}\right) .
\end{split}
\end{equation*}
Then
\begin{equation*}
\begin{split}
J_2(a) &= \sum_{m=1}^\infty  \int\limits_{ \varpi^{-m}  \mathcal{O}_F^\times}   \omega_\psi( (0, 0, z))  \Phi^0(a)       \tilde{f}_{\mathcal{W}(\tau,\psi_2), s} ^0  \left( \left(\begin{smallmatrix} a &\\&1  \end{smallmatrix}\right); \gamma u(0, 0, z)  \right)     
  dz     \\
 & = \sum_{m=1}^\infty  \int\limits_{ \varpi^{-m}  \mathcal{O}_F^\times}    \tilde{f}_{\mathcal{W}(\tau,\psi_2), s} ^0 \left( \left(\begin{smallmatrix} a & \\&1 \end{smallmatrix}\right);  \left(\begin{smallmatrix} 1 &&&\\&z^{-1}&-1&\\&&z&\\ &&&1\end{smallmatrix}\right)    \left(\begin{smallmatrix} 1 &&&\\&&1&\\&-1&&\\ &&&1\end{smallmatrix}\right)   \left(\begin{smallmatrix} 1 &&&\\&1&-z^{-1}&\\&&1&\\ &&&1\end{smallmatrix}\right)      \right)   \psi(z)dz \\
 & = \sum_{m=1}^\infty  \int\limits_{ \varpi^{-m}  \mathcal{O}_F^\times}    \tilde{f}_{\mathcal{W}(\tau,\psi_2), s} ^0  \left(\left(\begin{smallmatrix} a & \\&1 \end{smallmatrix}\right)  ; \left(\begin{smallmatrix} 1&&&\\&z^{-1}&&\\&&z&\\ &&&1\end{smallmatrix}\right)    \right)        \psi(z)dz \\
 &=   \sum_{m=1}^\infty  \int\limits_{ \varpi^{-m} \mathcal{O}_F^\times}   \gamma_\psi(z^{-1})^{-1} |z^{-1}|^{s+3/2} W_\tau^0  \left( \begin{smallmatrix} a & \\ & z^{-1} \end{smallmatrix}\right) \psi(z)dz.
\end{split}
\end{equation*}
Note that
\begin{equation*}
\begin{split}
W_\tau^0  \left( \begin{smallmatrix} a & \\ & z^{-1} \end{smallmatrix}\right) &=  W_\tau^0  \left(   \left( \begin{smallmatrix} z^{-1} & \\ & z^{-1} \end{smallmatrix}\right)  \left( \begin{smallmatrix} az & \\ & 1 \end{smallmatrix}\right)    \right) \\
&=\chi_1(z^{-1})\chi_2(z^{-1}) W_\tau^0  \left( \begin{smallmatrix} az & \\ & 1 \end{smallmatrix}\right)\\
  &=  \begin{cases}
\chi_1(z^{-1})\chi_2(z^{-1}) |az|^{\frac{1}{2}}  \cdot \frac{\chi_1(\varpi)^{\ord(az)+1}-\chi_2(\varpi)^{\ord(az)+1}}{\chi_1(\varpi)-\chi_2(\varpi)}       & \text{ if }\ord(z)\ge -\ord(a) \\
0 & \text{ if }\ord(z)<-\ord(a)
\end{cases}
\end{split}
\end{equation*}
where the last equality follows from Lemma \ref{lemma-CasselmanShalika}. If $\ord(a)=0$, then for any $z\in  \varpi^{-m} \mathcal{O}_F^\times$ with $m\ge 1$, we have $\ord(z)< -\ord(a)$, and hence $J_2(a)=0$. If $\ord(a)>0$, then
\begin{equation}
\begin{split}
J_2(a) = \sum_{m=1}^{\ord(a)}  \int\limits_{ \varpi^{-m} \mathcal{O}_F^\times}    \gamma_\psi(z^{-1})^{-1} |z^{-1}|^{s+3/2} W_\tau^0  \left( \begin{smallmatrix} a & \\ & z^{-1} \end{smallmatrix}\right) \psi(z)dz.
\end{split}
\label{eq-unramified-J2-reduction1}
\end{equation}


Now we assume $\ord(a)>0$. Then 
\begin{equation*}
\begin{split}
J_2(a)  &= \sum_{m=1}^{\ord(a)}  \int\limits_{ \varpi^{-m} \mathcal{O}_F^\times}   \gamma_\psi(z^{-1})^{-1} |z^{-1}|^{s+3/2}  \chi_1(z^{-1})\chi_2(z^{-1}) |az|^{\frac{1}{2}}  \cdot \frac{\chi_1(\varpi)^{\ord(az)+1}-\chi_2(\varpi)^{\ord(az)+1}}{\chi_1(\varpi)-\chi_2(\varpi)}      \psi(z)dz \\
&=    |a|^{\frac{1}{2}}  \sum_{m=1}^{\ord(a)}     \chi_1(\varpi^m)\chi_2(\varpi^m)  
\frac{\chi_1(\varpi)^{\ord(a)-m+1}-\chi_2(\varpi)^{\ord(a)-m+1}}{\chi_1(\varpi)-\chi_2(\varpi)}   \cdot q^{-m(s+1)}      \int\limits_{ \varpi^{-m} \mathcal{O}_F^\times} \gamma_\psi(z^{-1})^{-1}  \psi(z)dz \\
&=    |a|^{\frac{1}{2}}  \sum_{m=1}^{\ord(a)}     \chi_1(\varpi^m)\chi_2(\varpi^m)  
\frac{\chi_1(\varpi)^{\ord(a)-m+1}-\chi_2(\varpi)^{\ord(a)-m+1}}{\chi_1(\varpi)-\chi_2(\varpi)}   \cdot q^{-ms}      \int\limits_{\mathcal{O}_F^\times} \gamma_\psi( \varpi^{m} u^{-1})^{-1}  \psi( \varpi^{-m} u)du.
\end{split}
\end{equation*}
We remind the reader that here both $dz$ and $du$ are additively invariant. 

In Lemma \ref{lemma-local-integral-gamma-psi-vanishing} below, we will prove that only the term corresponding to $m=1$ contributes to the sum. Assuming Lemma \ref{lemma-local-integral-gamma-psi-vanishing} for the moment, then for $\ord(a)>0$ we have
\begin{equation*}
J_2(a)= |a|^{\frac{1}{2}} \cdot  \chi_1(\varpi)\chi_2(\varpi)  
\frac{\chi_1(\varpi)^{\ord(a)}-\chi_2(\varpi)^{\ord(a)}}{\chi_1(\varpi)-\chi_2(\varpi)}   \cdot q^{-s}      \int\limits_{\mathcal{O}_F^\times} \gamma_\psi( \varpi u^{-1})^{-1}  \psi( \varpi^{-1} u)du.
\end{equation*}
It follows from the formula $J_1(\mathbb{F},\psi,\chi^0)=q^{-\frac{1}{2}}$ in the proof of \cite[Lemma 1.12]{Szpruch2009} that
\begin{equation*}
\begin{split}
\int\limits_{\mathcal{O}_F^\times} \gamma_\psi( \varpi u^{-1})^{-1}  \psi(\varpi^{-1} u )  du =q^{-1/2} .
\end{split}
\end{equation*}
Therefore, 
\begin{equation*}
\begin{split}
J_2(a) &=   |a|^{1/2}\cdot   \frac{\chi_1(\varpi)^{\ord(a)}  - \chi_2(\varpi)^{\ord(a)} }{\chi_1(\varpi)-\chi_2(\varpi)} \cdot  \chi_1(\varpi)\chi_2(\varpi)  \cdot  q^{-s-1/2} .
\end{split}
\end{equation*}
Note that the above formula also applies to $J_2(a)=0$ when $\ord(a)=0$.

Thus, we will finish the proof of Proposition \ref{prop-unramified-J2} once we take care of Lemma \ref{lemma-local-integral-gamma-psi-vanishing}.
\end{proof}

\begin{thebibliography}{99}
\bibitem[Cassel]{CasselmanShalika1980II} W. Casselman and J. Shalika, \textit{The unramified principal series of {$p$}-adic groups. {II}. {T}he {W}hittaker function}, Compositio Math., \textbf{41}, (1980), no.~2, 207--231.
\bibitem[Szpruc]{Szpruch2009} D. Szpruch, \textit{Computation of the local coefficients for principal series representations of the metaplectic double cover of {${\rm SL}_2(\mathbb{F})$}}, J. Number Theory, \textbf{129} (2009), no.~9, 2180--2213.
\end{thebibliography}
\end{document}
